\chapter{Couche 9 — La chaîne DevSecOps}

Toutes les mesures précédentes se dégradent : une dépendance vieillit, un
développeur pressé affaiblit une règle, un secret se glisse dans un commit. La
chaîne d'intégration est ce qui les maintient vraies dans le temps.

\section*{Onze travaux, à chaque envoi}

\begin{center}
\small
\begin{tabularx}{\textwidth}{@{}p{3.2cm}p{2.5cm}X@{}}
\toprule
\textbf{Travail} & \textbf{Outil} & \textbf{Ce qu'il empêche d'entrer} \\
\midrule
\cmd{lint-python} & ruff & Du code mort, des imports inutilisés, des motifs
dangereux (\cmd{shell=True}, \cmd{assert} en production). Vérifie aussi que les
schémas du \cmd{README} sont à jour. \\
\cmd{tests-python} & pytest & Une régression fonctionnelle. 169 tests, dont le
cloisonnement entre comptes, le CSRF et l'anti-force-brute. \\
\cmd{sast-python} & bandit & Une faiblesse \emph{dans notre code} : secret en
dur, aléa non cryptographique, appel système non maîtrisé. \\
\cmd{sca-dependances} & pip-audit & Une dépendance portant une vulnérabilité
connue. C'est ainsi qu'a été repérée une faille de Werkzeug, corrigée en
passant de 3.1.8 à 3.1.9. \\
\cmd{secrets} & gitleaks & Un secret commis, y compris s'il a été retiré
depuis : l'outil parcourt \textbf{tout l'historique}. \\
\cmd{dast} & nuclei + \cmd{dast.sh} & Une faiblesse visible de l'extérieur sur
l'application \textbf{en train de tourner} : en-tête manquant, cookie mal
posé, route accessible sans session. \\
\cmd{lint-ansible} & ansible-lint & Un playbook non idempotent, une tâche sans
nom, un module obsolète. Profil \cmd{production}, le plus strict. \\
\cmd{images-docker} & hadolint + Trivy & Un \cmd{Dockerfile} mal écrit, et des
paquets vulnérables \emph{dans l'image livrée} --- ce que l'audit des
dépendances ne voit pas. \\
\cmd{integration} & \cmd{integration.sh} & Une plateforme qui ne démarre pas
sur une machine vierge. Les six scénarios du sujet. \\
\cmd{publier-images} & Syft + ghcr.io & Une image sans \textbf{inventaire de
composants} (SBOM) : on doit pouvoir répondre « cette image contient-elle la
bibliothèque vulnérable annoncée ce matin ? ». \\
\cmd{deployer} & Ansible & Un déploiement manuel, non reproductible. \\
\bottomrule
\end{tabularx}
\end{center}

\begin{notion}{SAST, SCA, DAST : trois angles, pas trois noms}
\textbf{SAST} lit \emph{notre} code sans l'exécuter --- il voit les erreurs que
nous écrivons. \textbf{SCA} confronte \emph{les dépendances} à une base de
vulnérabilités --- il voit les erreurs des autres, dont nous héritons.
\textbf{DAST} attaque l'application \emph{en fonctionnement} --- il voit ce
qu'aucune lecture de code ne montre : un en-tête oublié par la configuration
nginx, un cookie mal posé en production.

Les trois sont nécessaires parce qu'aucun ne voit ce que les deux autres
voient. La plupart des failles réelles tiennent à l'assemblage, pas à une
ligne.
\end{notion}

% -----------------------------------------------------------------------------
\mesure{Le seuil de blocage}{mesure 32 · \fichier{ci.yml}}

\menace
Une chaîne qui signale tout sans rien bloquer devient un bruit de fond qu'on
cesse de lire. Une chaîne qui bloque sur tout devient un obstacle qu'on
contourne.

\parade
Un seuil explicite, justifié, différent selon l'outil.

\danslecode
\cmd{hadolint} échoue au niveau \cmd{warning}, avec trois règles écartées et
chacune commentée --- par exemple l'épinglage des versions de paquets, écarté
parce que les dépôts Debian et Alpine ne conservent pas les anciennes versions :
une image épinglée cesserait de se construire à la première mise à jour du
dépôt, et c'est la version de la \emph{distribution} qui fixe déjà l'ensemble.

Trivy ignore les vulnérabilités \cmd{-{}-ignore-unfixed} : bloquer sur une faille
sans correctif disponible n'apporte rien --- on ne peut rien en faire --- et
noierait les failles réellement corrigeables.

\begin{retenir}
Un seuil bien choisi se reconnaît à ceci : chaque exception est écrite, datée
et justifiée dans un fichier versionné. Une exception non écrite est un oubli
qui se transmet.
\end{retenir}

\limite
Les seuils reflètent un jugement, qui vieillit. Une règle écartée aujourd'hui
pour de bonnes raisons peut devenir importante ; rien ne force à les relire.

% -----------------------------------------------------------------------------
\mesure{La version de Python, imposée}{mesure 33 · \fichier{ci.yml}}

\menace
Le projet passe les tests sur le poste du développeur et échoue au
déploiement, parce que la cible n'a pas la même version du langage. C'est
l'erreur la plus banale et la plus coûteuse, et elle a frappé ce projet : une
bibliothèque épinglée n'existait pas pour Python 3.10.

\parade
La chaîne force \cmd{PYTHON\_VERSION: 3.10} --- la version qu'Ubuntu 22.04
fournit, c'est-à-dire la cible réelle. Les tests tournent avec le Python de la
production, pas celui du développeur.

\limite
L'alignement porte sur la version du langage, pas sur l'ensemble du système :
seul le travail d'intégration, qui démarre réellement la pile, couvre le reste.

% -----------------------------------------------------------------------------
\mesure{Ce que la chaîne a réellement trouvé}{mesure 34 · bilan}

Une chaîne d'intégration ne vaut pas par sa configuration mais par ses prises.
Huit défauts ont été trouvés par la chaîne ou par l'exécution réelle, que la
relecture n'avait pas vus. Les trois plus instructifs :

\textbf{Un Faucheur qui détruisait sa propre plateforme.} Il reconnaissait les
machines louées par leur préfixe de nom ; les services de la pile
s'appellent \cmd{insacloud-web}, \cmd{insacloud-db}. Correction : l'identité
par étiquette (mesure 22).

\textbf{Une réconciliation qui détruisait les machines des autres nœuds.} Sur
un parc où deux nœuds partagent un démon Docker, chacun voyait les machines de
l'autre comme orphelines. Correction : l'étiquette \cmd{insacloud.worker}, et
deux tests --- dont un qui vérifie qu'une machine \emph{sans} étiquette reste
nettoyable, pour ne pas abandonner celles créées par une version antérieure.

\textbf{Une erreur 500 une fois sur trois, au démarrage seulement.} Gunicorn
tourne avec \cmd{-{}-preload} : il importe l'application \emph{avant} de forker
ses quatre workers, et l'initialisation de la base ouvre alors une connexion
PostgreSQL qui se range dans le pool du processus maître. Les quatre fils en
héritaient et se servaient de la \textbf{même socket} simultanément. Mesuré
avant correction : trois échecs sur huit premières requêtes parallèles, à
chacun des trois redémarrages d'essai ; après correction, vingt-quatre sur
vingt-quatre.

\begin{piege}{Le défaut qui ne se voit qu'au premier regard}
Ce dernier cas mérite d'être retenu pour sa forme : le message d'erreur
(\cmd{connection in transaction status INTRANS}) n'avait \emph{aucun} rapport
avec la requête en cours, et le symptôme disparaissait tout seul après
quelques appels. Un défaut qui se répare seul est un défaut qu'on croit
imaginaire --- et c'est exactement celui qu'aurait rencontré un visiteur
ouvrant la page pour la première fois.

Le test de non-régression a été validé en le faisant d'abord \textbf{échouer},
correction désactivée. Un test qu'on n'a pas vu échouer ne prouve rien.
\end{piege}
